\documentclass[11pt]{article}

% ---------- Packages ----------
\usepackage[a4paper,margin=1in]{geometry}
\usepackage{amsmath,amssymb,amsthm,mathtools}
\usepackage{bbm}
\usepackage{hyperref}
\usepackage{enumitem}

% ---------- Theorem environments ----------
\newtheorem{theorem}{Theorem}[section]
\newtheorem{proposition}[theorem]{Proposition}
\newtheorem{lemma}[theorem]{Lemma}
\newtheorem{corollary}[theorem]{Corollary}

\theoremstyle{definition}
\newtheorem{definition}[theorem]{Definition}
\newtheorem{remark}[theorem]{Remark}

% ---------- Commands ----------
\newcommand{\Bcal}{\mathcal{B}}
\newcommand{\OL}{\mathcal{O}_h^{(L)}}
\newcommand{\PhiL}{\Phi_h^{(L)}}
\newcommand{\Lp}[1]{L^{#1}}
\newcommand{\RL}{R_L}
\newcommand{\dL}{\delta(L)}

% ---------- Title ----------
\title{Algebraic and Geometric Reconstruction:\\
\large A Bridge via Conditional Variance}
\author{Patrick S.\ Mahon}
\date{}

\begin{document}
\maketitle

\begin{abstract}
Paper~IV introduces two witnesses for finite-sample reconstruction: the
algebraic witness $\delta(L)$ (worst-case $\sigma$-algebra approximation error)
and the geometric witness $(\mu \otimes \mu)(R_L)$ (mass of unseparated pairs
under the lag-$L$ equivalence relation).
This note identifies the exact object mediating between them --- the
\emph{conditional variance of indicators} --- and establishes a two-sided
comparison theorem.
The easy direction requires no conditions.
The hard direction requires that large fibres of the delay map be internally
mixed, a condition implied by ergodicity but strictly weaker.
As a corollary, reconstruction is equivalent to divergence of the collision
entropy of the delay-vector distribution, giving an information-theoretic
characterisation of reconstruction not present in Papers~I--IV.
\end{abstract}

% ============================================================
\section{Setup}
% ============================================================

Let $(X, \Bcal, \mu, T)$ be a measure-preserving system and $h \in \Lp{\infty}(\mu)$.
For $L \geq 0$ define the \emph{observable $\sigma$-algebra at lag~$L$}:
\[
  \OL \;:=\; \sigma\bigl(h \circ T^0,\, h \circ T^1,\, \ldots,\, h \circ T^L\bigr)
\]
and the \emph{lag-$L$ delay map} $\PhiL : X \to \mathbb{R}^{L+1}$,
$\PhiL(x) = (h(x), h(Tx), \ldots, h(T^L x))$.

\begin{definition}[Algebraic witness]
The \emph{$\sigma$-algebra approximation error at lag~$L$} is
\[
  \dL \;:=\; \sup_{S \in \Bcal} \inf_{E \in \OL} \mu(S \triangle E).
\]
\end{definition}

\begin{definition}[Equivalence relation and geometric witness]
Define the lag-$L$ equivalence relation $x \sim_L x'$ iff
$\PhiL(x) = \PhiL(x')$, and the corresponding unseparated-pairs set:
\[
  \RL \;:=\; \bigl\{(x,x') \in X \times X : x \sim_L x'\bigr\}.
\]
The \emph{geometric witness} is $(\mu \otimes \mu)(\RL)$.
\end{definition}

\begin{definition}[Fibre structure]
For $z \in \mathbb{R}^{L+1}$, let $F_z := (\PhiL)^{-1}(z)$ be the fibre over~$z$,
$p_z := \mu(F_z)$ the fibre mass, and $\mu_z$ the conditional measure on~$F_z$.
The push-forward $\nu_L := (\PhiL)_*\!\mu$ is the distribution of delay vectors.
\end{definition}

Three facts are immediate from the definitions:
\begin{enumerate}[label=(\alph*)]
  \item $\RL = \bigsqcup_z F_z \times F_z$, so
    $(\mu \otimes \mu)(\RL) = \int p_z^2\, d\nu_L(z) = \sum_z p_z^2$
    (in the discrete-fibre case).
  \item The conditional variance identity:
    for any $S \in \Bcal$,
    \[
      \mathbb{E}[\mathrm{Var}(\mathbf{1}_S \mid \OL)]
      \;=\; \int p_z \cdot \mu_z(S)\mu_z(S^c)\, d\nu_L(z).
    \]
  \item $\dL = \sup_S \mathbb{E}[\mathrm{Var}(\mathbf{1}_S \mid \OL)]$,
    since $\mathbb{E}[\mathrm{Var}(\mathbf{1}_S \mid \OL)]
    = \|\mathbf{1}_S - \mathbb{E}[\mathbf{1}_S \mid \OL]\|_{\Lp{2}(\mu)}^2$
    and $\dL = \sup_S \|\mathbf{1}_S - \mathbb{E}[\mathbf{1}_S \mid \OL]\|_{\Lp{2}(\mu)}^2$
    (Paper~IV, Lemma~3.1).
\end{enumerate}

% ============================================================
\section{The Bridge Identity}
% ============================================================

\begin{lemma}[Conditional variance identity]
\label{lem:bridge}
For any $S \in \Bcal$:
\[
  \mathbb{E}[\mathrm{Var}(\mathbf{1}_S \mid \OL)]
  \;=\;
  \frac{1}{2}
  \int_{\RL} |\mathbf{1}_S(x) - \mathbf{1}_S(x')|^2\,
  d(\mu \otimes \mu)(x, x').
\]
\end{lemma}

\begin{proof}
By disintegration, $(\mu \otimes \mu)|_{\RL} = \int \mu_z \otimes \mu_z\, d\nu_L(z)$.
On each fibre $F_z$:
\[
  \frac{1}{2}\int_{F_z \times F_z}
  |\mathbf{1}_S(x) - \mathbf{1}_S(x')|^2\, d(\mu_z \otimes \mu_z)
  \;=\; \mu_z(S)\mu_z(S^c).
\]
(Expand: $\int\int |\mathbf{1}_S(x) - \mathbf{1}_S(x')|^2 = 2[\mu_z(S)(1-\mu_z(S))]$.)
Integrating over fibres with weight $p_z = \mu(F_z)$:
\[
  \frac{1}{2}\int_{\RL}|\mathbf{1}_S(x) - \mathbf{1}_S(x')|^2\,d(\mu\otimes\mu)
  \;=\; \int p_z \cdot \mu_z(S)\mu_z(S^c)\,d\nu_L(z)
  \;=\; \mathbb{E}[\mathrm{Var}(\mathbf{1}_S \mid \OL)]. \qedhere
\]
\end{proof}

\begin{remark}
Lemma~\ref{lem:bridge} makes precise why the conditional variance is the
right mediating object:
its left-hand side is the quantity appearing in $\dL$ (algebraic),
and its right-hand side is an $S$-weighted version of $(\mu\otimes\mu)(\RL)$
(geometric).
Both measure the same failure of separation, but with different aggregation:
$\dL$ takes the worst case over $S$;
$(\mu\otimes\mu)(\RL)$ counts unseparated pairs with equal weight.
\end{remark}

% ============================================================
\section{The Two-Sided Comparison}
% ============================================================

\begin{theorem}[Easy direction --- unconditional]
\label{thm:easy}
\[
  \dL \;\leq\; \tfrac{1}{2}\,(\mu \otimes \mu)(\RL).
\]
\end{theorem}

\begin{proof}
From Lemma~\ref{lem:bridge} and $|\mathbf{1}_S(x) - \mathbf{1}_S(x')|^2 \leq 1$:
\[
  \mathbb{E}[\mathrm{Var}(\mathbf{1}_S \mid \OL)]
  \;\leq\; \tfrac{1}{2}\,(\mu \otimes \mu)(\RL)
\]
for every $S$.
Taking the supremum over $S$ gives $\dL \leq \tfrac{1}{2}(\mu\otimes\mu)(\RL)$.
\end{proof}

\begin{remark}
This says: large unseparated-pair mass forces large worst-case error.
No dynamical conditions are needed.
The constant $\tfrac{1}{2}$ is sharp: take $\mu$ uniform on two points,
$T = \mathrm{id}$, $h$ constant. Then $\RL = X \times X$,
$(\mu\otimes\mu)(\RL) = 1$, and $\dL = \tfrac{1}{2}$.
\end{remark}

For the reverse direction, the identity in Lemma~\ref{lem:bridge} shows that
the gap between $(\mu\otimes\mu)(\RL) = \int p_z^2\,d\nu_L(z)$ and
$\dL = \sup_S \int p_z \mu_z(S)\mu_z(S^c)\,d\nu_L(z)$ is controlled
by whether large fibres ($p_z$ large) are internally mixed.

\begin{definition}[Fibre mixing condition]
\label{def:fibre-mixing}
Say that $S^* \in \Bcal$ is \emph{$c$-mixing across fibres} if there exists
$c > 0$ such that
\[
  \mu_z(S^*)\mu_z((S^*)^c) \;\geq\; c \cdot p_z
  \quad \text{for $\nu_L$-a.e.\ } z \text{ with } p_z > 0.
\]
\end{definition}

\begin{theorem}[Hard direction --- conditional]
\label{thm:hard}
If the event $S^*$ achieving $\dL$ is $c$-mixing across fibres, then
\[
  (\mu \otimes \mu)(\RL) \;\leq\; \frac{1}{c}\,\dL.
\]
\end{theorem}

\begin{proof}
From the fibre mixing condition:
\[
  \dL
  \;\geq\; \int p_z \cdot \mu_z(S^*)\mu_z((S^*)^c)\,d\nu_L(z)
  \;\geq\; c \int p_z^2\,d\nu_L(z)
  \;=\; c\,(\mu\otimes\mu)(\RL). \qedhere
\]
\end{proof}

\begin{remark}[Fibre mixing vs ergodicity]
\label{rem:ergodicity}
The fibre mixing condition (Definition~\ref{def:fibre-mixing}) is strictly
weaker than ergodicity.
It requires only that the adversarial event $S^*$ intersects every positive-mass
fibre in a non-trivial fraction.
Ergodicity implies this: under an ergodic $T$, if $S^*$ achieves $\dL > 0$
then $S^*$ is not $\OL$-measurable, so $S^*$ must straddle some fibres, and
by ergodicity the mass of non-straddled fibres is zero.
But ergodicity is not necessary: the condition fails only when $S^*$ is
``aligned'' with the fibre structure in the sense that large fibres are
entirely inside or outside $S^*$.
\end{remark}

\begin{corollary}[Equivalence under fibre mixing]
\label{cor:equiv}
Under the fibre mixing condition:
\[
  c\,(\mu\otimes\mu)(\RL) \;\leq\; \dL \;\leq\; \tfrac{1}{2}\,(\mu\otimes\mu)(\RL).
\]
In particular:
\[
  \dL \to 0 \quad\iff\quad (\mu\otimes\mu)(\RL) \to 0.
\]
\end{corollary}

\begin{proof}
Theorems~\ref{thm:easy} and~\ref{thm:hard} combined.
\end{proof}

% ============================================================
\section{Entropy Characterisation of Reconstruction}
% ============================================================

The geometric witness $(\mu\otimes\mu)(\RL) = \int p_z^2\,d\nu_L(z)$ is the
collision probability of the delay-vector distribution $\nu_L$.
Its logarithm is the (negative) \emph{Rényi-$2$ entropy} (collision entropy) of $\nu_L$:
\[
  H_2(\nu_L) \;:=\; -\log\bigl((\mu\otimes\mu)(\RL)\bigr)
  \;=\; -\log\!\int p_z^2\,d\nu_L(z).
\]

\begin{corollary}[Entropy characterisation of reconstruction]
\label{cor:entropy}
Under the fibre mixing condition of Definition~\ref{def:fibre-mixing}:
\[
  \dL \to 0
  \quad\iff\quad
  H_2(\nu_L) \to +\infty.
\]
That is, reconstruction holds if and only if the collision entropy of
the delay-vector distribution diverges.
\end{corollary}

\begin{proof}
$\dL \to 0 \iff (\mu\otimes\mu)(\RL) \to 0$ (Corollary~\ref{cor:equiv})
$\iff e^{-H_2(\nu_L)} \to 0 \iff H_2(\nu_L) \to +\infty$.
\end{proof}

\begin{remark}[Interpretation]
Corollary~\ref{cor:entropy} gives a purely information-theoretic
characterisation of reconstruction: the observable $h$ reconstructs the system
$(X, \Bcal, \mu, T)$ at lag $L$ if and only if the distribution of the delay
vector $\PhiL(x)$ becomes increasingly spread --- in the sense of Rényi-$2$
entropy --- as $L$ grows.

This characterisation is not present in Papers~I--IV.
It does not require:
\begin{itemize}
  \item a metric on $X$,
  \item smoothness of $T$ or $h$,
  \item a mixing rate,
  \item or the Takens embedding conditions.
\end{itemize}
The collision entropy $H_2(\nu_L)$ is estimable from data: it is the
negative log of $\sum_z (\text{fraction of sample in fibre } z)^2$,
a standard diversity index computable from the delay-vector histogram.
Under Paper~IV's standing hypotheses (compact smooth $X$, smooth positive $\mu$,
$T \in C^r$), the fibre mixing condition holds under ergodicity, and the
entropy characterisation then provides a third, data-driven witness for
reconstruction complementing $\delta(L)$ and the delay-map separation.
\end{remark}

\begin{remark}[Relation to Paper~IV witnesses]
The three witnesses now in scope are:
\begin{enumerate}[label=(\roman*)]
  \item $\hat\delta(L,n)$ --- the algebraic witness (Paper~IV, \S5);
  \item $\hat d_L(x,x')$ --- the delay-map separation (Paper~IV, \S6);
  \item $\hat H_2(\nu_L^{(n)}) := -\log\bigl(\frac{1}{n^2}\sum_{i,j}
        \mathbf{1}[\PhiL(x_i) = \PhiL(x_j)]\bigr)$ --- the empirical
        collision entropy.
\end{enumerate}
All three are computable from the time series alone.
Corollaries~\ref{cor:equiv} and~\ref{cor:entropy} show that (i) and (iii)
converge to zero (resp.\ diverge) together under fibre mixing,
making them asymptotically equivalent witnesses for reconstruction.
Witness~(ii) provides the geometric certificate that the state space is
being faithfully separated.
\end{remark}

\end{document}
